% =====================================================================
%  Основное содержимое отчёта.
% =====================================================================

\labpart{Краткое описание разработанного приложения}

\texttt{SecureApp} -- клиент-серверное веб-приложение, состоящее из двух
независимо развёртываемых частей.

\begin{itemize}
  \item \texttt{Backend/SecureApp.Api} -- REST API на ASP.NET Core (.NET),
  хранящее данные в локальной базе SQLite через Entity Framework Core.
  Реализует все операции задания: регистрацию (\texttt{POST
  /api/auth/register}), вход (\texttt{POST /api/auth/login}), выход
  (\texttt{POST /api/auth/logout}) и по четыре операции CRUD и поиска для
  каждого из двух типов данных (\texttt{/api/confidential},
  \texttt{/api/public}).
  \item \texttt{Frontend/secure-app-ui} -- одностраничное приложение на
  React (TypeScript) с компонентами Material UI: формы входа и
  регистрации, а также панель с двумя вкладками -- «Confidential data» и
  «Public data», -- в каждой из которых таблица записей, строка поиска и
  диалог создания/редактирования.
\end{itemize}

Каждая запись (конфиденциальная и неконфиденциальная) принадлежит ровно
одному пользователю: сервер связывает её с \texttt{OwnerId} из проверенного
токена сессии, а не с значением, присланным клиентом, поэтому чтение,
изменение, удаление и поиск всегда выполняются в границах данных текущего
пользователя.

\labpartbreak{Анализ защищённости приложения}

\subsection*{1. Политика безопасности операционной системы}

Приложение не хранит секреты в коде и не рассчитывает на права
операционной системы для защиты данных -- оно рассчитывает на собственный
уровень контроля доступа (см. далее), что делает его защищённость
переносимой между ОС. Тем не менее ОС-уровень остаётся значимым: ключ
подписи JWT и ключ шифрования конфиденциальных данных вынесены из кода в
конфигурацию (\texttt{appsettings.Development.json} -- только для
локальной разработки) с явным указанием, что в эксплуатации они должны
задаваться переменными окружения (\texttt{Jwt\_\_Key},
\texttt{Encryption\_\_Key}) или менеджером секретов, а не файлом,
попадающим в систему контроля версий; в этом случае получить их может
только тот, у кого есть права на переменные окружения процесса, то есть
тот же пользователь ОС, что и владелец процесса сервера. Файл базы данных
SQLite создаётся с правами по умолчанию текущего пользователя ОС
(\texttt{umask} процесса), то есть недоступен на чтение другим локальным
пользователям без дополнительных привилегий.

\subsection*{2. Политика безопасности компонентов приложения и их взаимодействия}

Взаимодействие «клиент -- сервер» построено на нескольких независимых
рубежах защиты, которые не дублируют, а дополняют друг друга.

Во-первых, аутентификация не хранит пароль -- при регистрации пароль
преобразуется алгоритмом PBKDF2-HMACSHA256 со случайной солью на
пользователя и 210~000 итераций (см. \texttt{PasswordHasher.cs}), после
чего сравнивается только результат хеширования. Во-вторых, состояние
сессии передаётся не в теле ответа и не в заголовке, а в
cookie-заголовке с атрибутами \texttt{HttpOnly} (недоступен
JavaScript-коду страницы -- защита от кражи токена через XSS),
\texttt{Secure} (передаётся только по защищённому соединению) и
\texttt{SameSite=Strict} (не отправляется при переходах с посторонних
сайтов -- защита от CSRF), что и подтверждено в
\autoref{lst:secureapp-output} (Шаг 3). В-третьих, каждый защищённый
маршрут API отмечен атрибутом \texttt{[Authorize]} и дополнительно
фильтрует данные по идентификатору владельца, извлечённому из подписанного
токена, а не переданному клиентом явно (см.
\autoref{lst:secureapp-owner-scope}), поэтому компонент хранения данных и
компонент разграничения доступа физически не могут разойтись -- условие
владения встроено в каждый запрос к базе.

\begin{codelisting}{csharp}{Каждый запрос к конфиденциальным данным ограничен владельцем записи}{secureapp-owner-scope}
private int CurrentUserId => int.Parse(User.FindFirstValue(ClaimTypes.NameIdentifier)!);

[HttpGet("{id:int}")]
public async Task<ActionResult<RecordResponse>> Get(int id)
{
    var record = await db.ConfidentialRecords
        .SingleOrDefaultAsync(r => r.Id == id && r.OwnerId == CurrentUserId);
    return record is null ? NotFound() : Ok(ToResponse(record));
}
\end{codelisting}

Обращение постороннего пользователя к чужой записи намеренно возвращает
\texttt{404 Not Found}, а не \texttt{403 Forbidden}: с точки зрения
пользователя, не владеющего записью, между «записи не существует» и
«записи существует, но доступ запрещён» не должно быть заметной разницы --
иначе сам код ответа стал бы каналом утечки информации о существовании
чужих идентификаторов.

Наконец, попытки входа с неверным паролем не различаются с попытками входа
под несуществующим именем пользователя -- в обоих случаях возвращается
одно и то же сообщение (\autoref{lst:secureapp-output}, Шаг 4), что не
позволяет использовать форму входа для перебора существующих в системе
имён пользователей.

\subsection*{3. Сетевая политика безопасности}

Взаимодействие фронтенда и API ограничено списком разрешённых источников
(CORS), настроенным по одному конкретному адресу фронтенда, а не по маске
или символу \texttt{*}, причём заголовок
\texttt{Access-Control-Allow-Credentials} выставляется только для этого
источника -- запрос с постороннего источника не получает в ответе заголовка
\texttt{Access-Control-Allow-Origin} вовсе (\autoref{lst:secureapp-output},
Шаг 10), то есть браузер блокирует его на своей стороне ещё до того, как
ответ окажется доступен вызывающему скрипту. Дополнительно ко всем ответам
добавляются заголовки \texttt{X-Content-Type-Options: nosniff},
\texttt{X-Frame-Options: DENY} и \texttt{Referrer-Policy: no-referrer}
(см. \texttt{Program.cs}), снижающие, соответственно, риск MIME-sniffing
атак, встраивания страницы в чужой \texttt{<iframe>} (clickjacking) и
утечки внутренних путей приложения через заголовок \texttt{Referer} при
переходе по внешним ссылкам. Транспорт защищён на уровне промежуточного ПО
\texttt{UseHttpsRedirection}, форсирующего HTTPS в развёрнутой
конфигурации. Наконец, на сетевом канале входа реализована защита от
перебора паролей: после пяти неудачных попыток входа под одним и тем же
именем дальнейшие попытки отклоняются кодом \texttt{429 Too Many Requests}
на 30 секунд вне зависимости от корректности пароля -- продемонстрировано
в \autoref{lst:secureapp-output} (Шаг 5), где шестая попытка отклонена,
хотя пароль в ней уже верен.

\labpartbreak{Анализ защищённости конфиденциальной информации}

Различие между конфиденциальными и неконфиденциальными данными в
\texttt{SecureApp} проведено не только на уровне доступа (см. выше), но и
на уровне хранения. Неконфиденциальные записи хранятся в столбцах типа
\texttt{TEXT} открытым текстом. Конфиденциальные записи перед записью в
базу шифруются алгоритмом AES-256-GCM с ключом, производным от
конфигурационного секрета (\autoref{lst:secureapp-crypto}): используется
аутентифицированное шифрование, то есть попытка подменить
зашифрованные байты в базе напрямую (в обход API) приведёт не к
подстановке произвольных данных, а к ошибке проверки тега аутентификации
при расшифровке.

\begin{codelisting}[fontsize=\footnotesize]{csharp}{Шифрование конфиденциального поля перед сохранением (AES-256-GCM)}{secureapp-crypto}
public byte[] Encrypt(string plaintext)
{
    var nonce = RandomNumberGenerator.GetBytes(NonceSize);
    var plainBytes = Encoding.UTF8.GetBytes(plaintext);
    var cipherBytes = new byte[plainBytes.Length];
    var tag = new byte[TagSize];

    using var aes = new AesGcm(_key, TagSize);
    aes.Encrypt(nonce, plainBytes, cipherBytes, tag);
    // Результат: [nonce][ciphertext][tag] -- всё, что попадает в столбец BLOB.
    ...
}
\end{codelisting}

Утверждение о шифровании проверено не по коду, а экспериментально: в базу
была записана конфиденциальная запись со значением
\texttt{4111-0000-0000-0000} и неконфиденциальная запись со значением
\texttt{Office closed Friday}, после чего необработанный файл базы данных
(включая WAL-журнал SQLite) был просканирован утилитой \texttt{strings}.
Конфиденциальное значение не встретилось в файле ни в каком виде, тогда
как неконфиденциальное было найдено открытым текстом
(\autoref{lst:secureapp-output}, Шаг 11) -- то есть даже полная утечка
файла базы данных (например, копированием диска) не раскрывает
конфиденциальные поля без отдельно хранимого ключа шифрования, тогда как
неконфиденциальные поля таким сценарием раскрываются полностью, что
ожидаемо и соответствует их статусу.

Полный протокол демонстрации, включающий регистрацию, вход, разграничение
данных между двумя пользователями, поиск, CORS и завершение сессии,
приведён ниже.

\codelistingfile[fontsize=\scriptsize]{text}{lab/code/output.txt}{Реальный протокол проверки SecureApp (скрипт \texttt{verify.sh})}{secureapp-output}
